\section{Preliminaries in nonlinear potential theory}
\label{sec:2}
We first remind that the operator $\Delta_p$ is defined by $\Delta_p f = \div \big(\abs{\nabla f}^{p-2}\nabla f\big)$, for $f \in \CS^2$ with nonvanishing gradient.
As far as basic principles and regularity for \emph{weakly} $p$-harmonic functions are concerned, we just refer the reader to
\cite{dibenedetto_alphalocalregularityweak_1983,evans_newprooflocal1_1982,heinonen_asuperharmonicfunctionssupersolutionsdegenerate_1988,ladyzenskaja_linearquasilinearellipticequations_1968,lewis_regularityderivativessolutionscertain_1983,lieberman_boundaryregularitysolutionsdegenerate_1988,serrin_localbehaviorsolutionsquasilinear_1964,tolksdorf_dirichletproblemquasilinearequations_1983,trudinger_harnacktypeinequalitiestheir_1967,uraltseva_degeneratequasilinearellipticsystems_1968,valtorta_laplaceoperatorriemannianmanifolds_2013}, and~\cite{holopainen_nonlinearpotentialtheoryquasiregular_1990,holopainen_volumegrowthgreenfunctions_1999} for the theory of $p$-capacitary potentials in exterior domains (see also~\cite[Chapter 1]{benatti_monotonicityformulasnonlinearpotential_2022} and reference therein), including existence issues. The function $w_p$ that solves~\eqref{eq:pb-intro} can be written as $w_p= -(p-1)\log u_p$, where $u_p \in \CS^{1,\beta}_{\loc}(M \smallsetminus \Int \Omega)$ is the solution to
\begin{equation}\label{eq:p-cap-potential}
 \begin{cases}
 \Delta^{\tp}_g u_p = 0 &\text{on $M \smallsetminus \Int \Omega$,}\\
 u_p=1 & \text{on $\partial \Omega$,}\\
 u_p\to 0 & \text{as $\dist(x,\partial \Omega)\to +\infty$,}
 \end{cases}
\end{equation}
the following definition of strongly $p$-nonparabolic Riemannian manifold is consistent with that of strong nonparabolicity~\cite{ni_meanvaluetheoremsmanifolds_2007} and with the limit case of strong $1$-nonparabolicity~\cite{benatti_isoperimetricriemannianpenroseinequality_2022}.

\begin{Definition}[strongly $p$-nonparabolic]
We say that a $3$-dimensional Riemannian manifold~$(M,g)$ with compact possibly empty boundary is \emph{strongly $p$-nonparabolic}, $1<p<3$, if there exists a solution to~\eqref{eq:pb-intro} for some $\Omega \subseteq M$ closed bounded with smooth boundary homologous to $\partial M$.
\end{Definition}

\begin{Remark}
By the maximum principle, in a strongly $p$-nonparabolic manifold every $\Omega$ with $\CS^{1}$-boundary homologous to $\partial M$ admits a solution to~\eqref{eq:pb-intro}.
\end{Remark}

This definition naturally comes with the notion of the $p$-capacity of a compact subset $K \subset M$, which we recall to be
\begin{equation*}%\label{eq:p-capacity}
 \ncapa_p(K) = \inf \set{\frac{1}{4\pi}\bigg(\frac{p-1}{3-p}\bigg)^{p-1}\int_{M\smallsetminus K} \abs{\D v}^p \dif \mu\st v \in \CS_c^\infty(M),\, v \geq 1\text{ on } K}.
\end{equation*}
When the boundary of $K$ is sufficiently regular, the above infimum is realized by the function~$u_p$ solving~\eqref{eq:p-cap-potential}. A fundamental property of the $p$-capacity is that it is monotone with respect to the standard inclusion of sets. More specifically, if we consider sublevel sets $\Omega_t=\set{w_p \leq t}$ of solutions to~\eqref{eq:pb-intro}, we have that their $p$-capacities grow exponentially with respect to the arrival time parameter $t$. This is
completely analogous to the exponential growth of the area along the IMCF. We recall this useful property in the following lemma, proved in~\cite[Lemma~3.8]{holopainen_nonlinearpotentialtheoryquasiregular_1990}.

\begin{Lemma}
Let $(M,g)$ be a $3$-dimensional Riemannian manifold with compact possibly empty boundary $\partial M$. Let $\Omega\subseteq M$ be a closed bounded subset homologous to $\partial M$ with $\CS^1$-boundary, $w_p$ the solution to~\eqref{eq:pb-intro} starting at $\Omega$ and $\Omega_t= \set{w_p \leq t}$, $1<p<3$. We have
\[
 \ncapa_p(\partial \Omega_t) = \ee^{t} \ncapa_p(\partial \Omega)= \frac{1}{4\pi}\int_{\partial \Omega_t} \bigg(\frac{ \abs{ \D w_p}}{3-p}\bigg)^{p-1}\dif \sigma.
\]
\end{Lemma}
